% pics.tex - original TikZ drawings of the LUCE poster (Conclusions icons).
% The Methods cells and the big drawing of the model are in methods.tex.
% All colours come from palette.tex. Each macro takes the scale as argument.

\tikzset{
  ray/.style={line width=2.2pt, line cap=round, >={Stealth[length=9pt,width=8pt]}},
  thinray/.style={line width=1.4pt, line cap=round, >={Stealth[length=7pt,width=6pt]}},
  sea/.style={fill=lucereal!14, draw=none},
  lab/.style={font=\fpic, text=lucetext, inner sep=2pt},
  sat/.pic={
    \fill[luceframe] (-0.28,-0.2) rectangle (0.28,0.2);
    \fill[luceframe!55] (-1.0,-0.13) rectangle (-0.36,0.13);
    \fill[luceframe!55] (0.36,-0.13) rectangle (1.0,0.13);
  }
}

% --- Rough surface: facet, reflected and refracted rays (Background 2, B) --
\newcommand{\picSurface}[1]{%
\begin{tikzpicture}[scale=#1]
  \fill[sea] plot[domain=-4.2:4.2,samples=60] (\x,{0.45*sin(75*\x)}) -- (4.2,-2.6) -- (-4.2,-2.6) -- cycle;
  \draw[lucereal, line width=2pt] plot[domain=-4.2:4.2,samples=60] (\x,{0.45*sin(75*\x)});
  % Facet at x = 0 (slope of the sine there: 0.45*75*pi/180 = 0.589).
  \draw[luceframe, line width=4pt, line cap=round] (-0.75,-0.44) -- (0.75,0.44);
  \draw[thinray, luceframe, ->] (0,0) -- (-1.02,1.72) node[lab, anchor=south] {$\hat{\mathbf{n}}$};
  \draw[ray, lucetime, ->] (-0.4,3.4) -- (0,0);
  \draw[ray, lucetime, ->] (0,0) -- (-2.9,1.95);
  \draw[ray, lucetime, dashed, ->] (0,0) -- (-0.25,-2.3);
  \node[lab, anchor=west] at (-0.2,3.0) {$\hat{\mathbf{k}}_\mathrm{in}$};
  \node[lab, anchor=south] at (-3.0,2.0) {$\hat{\mathbf{k}}_\mathrm{ref}$};
  \node[lab, anchor=west] at (-0.15,-1.9) {$\hat{\mathbf{k}}_\mathrm{tr}$};
  % Shadowed ray.
  \draw[thinray, lucelimit, dashed, ->] (4.2,1.25) -- (2.75,0.62);
  \node[lab, anchor=south, text=lucelimit] at (3.0,1.25) {shadowed};
\end{tikzpicture}}

% --- Real depth against time-based depth (Background 3, F) ----------------
\newcommand{\picPath}[1]{%
\begin{tikzpicture}[scale=#1]
  \fill[sea] (-4.2,-6.3) rectangle (4.2,0);
  \draw[lucereal, line width=2pt] (-4.2,0) -- (4.2,0);
  \draw[ray, lucetime] (-1.2,1.4) -- (-0.9,0);
  % Multiply scattered path.
  \draw[ray, lucereal] (-0.9,0) -- (-0.5,-1.3) -- (0.5,-2.0) -- (0.2,-2.9);
  \draw[ray, lucereal, ->] (0.2,-2.9) -- (-0.35,0);
  \foreach \p in {(-0.5,-1.3),(0.5,-2.0),(0.2,-2.9)} \fill[lucereal] \p circle (0.13);
  \draw[thinray, lucereal, dashed] (0.2,-2.9) -- (2.3,-2.9);
  \node[lab, anchor=west, text=lucereal] at (2.3,-2.9) {$z$};
  % Straight path of the same length, as the lidar assumes.
  \draw[ray, lucetime, dashed] (-2.2,0) -- (-1.6,-4.6);
  \fill[lucetime] (-1.6,-4.6) circle (0.13);
  \draw[thinray, lucetime, dashed] (-1.6,-4.6) -- (2.3,-4.6);
  \node[lab, anchor=west, text=lucetime] at (2.3,-4.6) {$z_t$};
  \node[lab] at (0.3,-5.55) {same time of flight};
\end{tikzpicture}}

% --- Local estimate: every vertex is scored toward the telescope (E) -------
\newcommand{\picLE}[1]{%
\begin{tikzpicture}[scale=#1]
  \fill[sea] (-4.2,-4.2) rectangle (4.2,0);
  \draw[lucereal, line width=2pt] (-4.2,0) -- (4.2,0);
  \path (-3.0,3.0) pic[rotate=-20,scale=0.9]{sat};
  \draw[ray, lucetime, ->] (-2.9,2.8) -- (-0.2,0);
  \draw[ray, lucereal] (-0.2,0) -- (0.6,-1.4) -- (1.9,-2.0) -- (1.5,-3.3) -- (2.9,-3.8);
  \foreach \x/\y in {0.6/-1.4,1.9/-2.0,1.5/-3.3} {
    \fill[lucereal] (\x,\y) circle (0.13);
    \draw[thinray, luceorder1, dashed, ->] (\x,\y) -- ({\x-0.16*(-\y)},0) -- ({\x-0.16*(-\y)-1.6},1.75);
  }
  \draw[thinray, luceorder2, dashed, ->] (-0.2,0) -- (-2.3,2.35);
  \node[lab, anchor=west, text=luceorder1] at (2.2,-1.2) {$E(z)$};
  \node[lab, anchor=east, text=luceorder2!80!black] at (-1.9,1.0) {$E_\mathrm{glint}$};
\end{tikzpicture}}
